\begin{theorem}[The explicit fourth-root block construction on the torus]%
    \label{thm:peps_torus_block_multiplicity_state}
    \lean{TNLean.PEPS.blockMatrixRepresentation,
        TNLean.PEPS.blockFourthRootWeight,
        TNLean.PEPS.blockFourthRootWeight_sq_mul,
        TNLean.PEPS.blockFourthRootWeight_commute,
        TNLean.PEPS.torusBondRegrouping_blockFourthRootWeight,
        TNLean.PEPS.torusBondNetwork_blockFourthRootWeight_norm}
    \leanok
    \uses{thm:peps_torus_weighted_blocks_state,
        def:peps_multiplicity_restored_representation,
        thm:peps_torus_multiplicity_bond_state,
        lem:peps_matching_sector_bond_inclusion,
        thm:peps_full_multiplicity_bond_map,
        thm:peps_multiplicity_bond_product_overlap}
    Let $D_i:G\to M_{d_i}(\mathbb C)$ be chosen matrix representations of a
    finite group, with $d_i>0$. Define the representation and weight
    \begin{align}
        U(g) & =\bigoplus_iD_i(g),            &
        W    & =\bigoplus_i d_i^{1/4}I_{d_i}.
    \end{align}
    The weight commutes with every $U(g)$, and
    \begin{align}
        W^2U(g)=\bigoplus_i\sqrt{d_i}\,D_i(g).
    \end{align}
    Let $H$ be the actual torus state obtained by dressing the normalized
    averaging site of $U$ with $W$ on each virtual leg. Put
    $L(g)=\bigoplus_iD_i(g)\otimes I_{d_i}$, and let $K_L$ be its actual
    normalized averaging-site torus state. If $\mathcal B$ groups the original
    physical legs into their oriented bonds, $E$ includes matching-sector
    bond labels, and $T$ restores multiplicities, then
    \begin{align}
        \mathcal B H               & \in\operatorname{range}(E^{\otimes 2N}), &
        T^{\otimes 2N}\mathcal B H & =\mathcal B K_L.
    \end{align}
    Here $N$ is the number of vertices, and the torus circumferences are positive.
    Moreover, the two states have equal squared norms,
    $\langle H,H\rangle=\langle K_L,K_L\rangle$.
    Thus the support and transformation of the actual contractions follow
    from the explicit block construction in \cite[Section~7]{Schuch2010PEPS}.
    This algebraic assertion does not assume that the chosen blocks enumerate
    all irreducibles. The existence of that enumeration and its unitary Fourier
    identification with the regular representation remain separate; see
    \cite{gap:rmp_peps_quantum_double_g_isometry}.
\end{theorem}

\begin{proof}\leanok
    Multiply the block-scalar matrices and use
    $(\sqrt{\sqrt{d_i}})^2=\sqrt{d_i}$. The scalar blocks commute with each
    representation block. Expand the actual dressed torus contraction into
    its coherent product of physical bond vectors. Every factor is the
    weighted block matrix displayed above, so matching-sector inclusion
    reconstructs it. The multiplicity-restoring bond map replaces each
    $\sqrt{d_i}D_i(g)$ by $D_i(g)\otimes I_{d_i}$.
    The coherent product-map identity then gives the actual torus contraction
    of $L$, with the same factor $|G|^{-N}$ on both sides.
    The product bond map preserves overlaps on its matching-sector support,
    and physical regrouping preserves all overlaps. Combining these facts
    with the derived support and state equality gives the norm identity.
\end{proof}
